% !TEX program = xelatex
\IfFileExists{main.tex}{\documentclass[main.tex]{subfiles}}{
  \IfFileExists{../../main.tex}{\documentclass[../../main.tex]{subfiles}}{
    \IfFileExists{../main.tex}{\documentclass[../main.tex]{subfiles}}{
      \documentclass[Book\space NL/main.tex]{subfiles}
    }
  }
}
\begin{document}
\ifSubfilesClassLoaded{
  \ifcsname appendixchapterparent\endcsname
  \else
    \standaloneappendixchaptersetup{C}{2}
  \fi
}{}
\chapter{Virology: Viruses and Viral Diseases}
\labch{appendix-c2-virology}
\margintoc

While the preceding sections of our text discuss medical microbiology from a systems-based perspective, this section introduces the foundations of medical virology and provides an overview of this ever-expanding field. This section is designed to present a ``microbe parade'' of sorts (following Murray's approach\footnote{Hopefully it is abundantly clear that we intend to focus on knowledge with greatest relevance in practice for the Southeast Asian practitioner.}) and act as a summary of the knowledge presented in the preceding sections.

Personally, I find that medical microbiology is one of the most difficult topics in all of clinical medicine, if not the most difficult, thanks in part to the massive body of knowledge that must be remembered and understood by medical students. As such, I hope that this section serves as a starting point for students looking to develop their understanding and more effectively tackle examinations involving virology and the pathophysiology of infection. Additionally, I note with thanks that the medical virology section of this text is based on the core concepts presented in Section 5 (``Virology,'' chapters 35--56) of \textit{Medical Microbiology}, 10th ed., by \textcite{Murray2026MedicalMicrobiology}.

Content coverage of this topic is as follows [quoted verbatim, sic. from specifications]:
\begin{itemize}
  \item B1.5.1 -- Principles of medical microbiology and parasitology
  \begin{enumerate}
    \item Human pathogen classification (bacteria, virus, fungus, parasite, prion)
    \item Basic Human microbiota
    \item Basic biosafety for general practitioner (Biosafety Level, Personal protective equipment, sterilization, disinfection)
    \item Microbes in medical biotechnology applications (probiotics, viral vector vaccine, gene therapy, oncolytic virus, phage therapy)
    \item Vectors control, and prevention (mosquito, flea, tick, mite and fly)
    \item Emerging infectious diseases (e.g., COVID 19)
    \item Principle of drug resistance pathogens
    \item Principle of serologic studies
    \item Principles of molecular diagnosis
  \end{enumerate}
  \item B1.5.3 -- General Principles of medical virology
  \begin{enumerate}
    \item Classification of Medically Important Viruses
    \item Structure, genome (general concept for understanding antivrial agents, vaccine)
    \item Replication (general concept for understanding antiviral agents, antibody)
    \item Pathogenisis
    \item Principles of cultivation, assay and laboratory diagnosis
    \item Zoonosis, vector-bourne disease of pathogenic virus: dengue hemorrhagic fever, zika, Japanese encephalitis, chikungunya, rabies
  \end{enumerate}
\end{itemize}

\begin{center}
\captionof{table}{Alignment of the virology chapter with NL content objectives. Source: Medical Council Notification 12/2555, Part 1A; coverage reflects the current content of this chapter.}
\labtab{virology_nl_content_alignment}
\footnotesize
\setlength{\tabcolsep}{2pt}
\renewcommand{\arraystretch}{1.12}
\begin{tabular}{@{}>{\RaggedRight\arraybackslash}p{0.14\textwidth}>{\RaggedRight\arraybackslash}p{0.44\textwidth}>{\RaggedRight\arraybackslash}p{0.39\textwidth}@{}}
\toprule
\textbf{NL code} & \textbf{Content objective} & \textbf{Coverage in this chapter} \\
\midrule
B1.5.7.1 & Microbial classification & Partial: viral genome grouping is introduced; broader microbial classification and formal virus taxonomy need further coverage. \\
B1.5.7.3 & Viruses and viral diseases & Partial: foundational virology and classification are introduced; detailed replication and clinical disease content remain to be covered. \\
B1.5.7.3.1 & Physical and chemical properties & Partial: capsid and envelope structure and environmental stability are introduced, but properties are not reviewed systematically. \\
B1.5.7.3.2 & Replication & Partial: host-cell dependence and Baltimore routes to mRNA are introduced; complete replication cycles remain to be covered. \\
B1.5.7.3.3 & Genetics & Partial: DNA/RNA genome types are introduced; viral genetics and genetic variation are not covered. \\
B1.5.7.3.4 & Principles of cultivation, assay, and laboratory diagnosis & Not yet covered. \\
B1.5.7.3.5 & Molecular basis of pathogenesis & Partial: immune evasion is mentioned, but mechanisms of pathogenesis are not explained. \\
B1.5.7.3.6 & Pathophysiology of infection & Not yet covered. \\
B1.5.7.3.7 & Latent and persistent infections & Not yet covered. \\
B1.5.7.3.8 & Epidemiology & Partial: general transmission routes are listed; virus-specific epidemiology is not covered. \\
B1.5.7.3.9 & Oncogenic viruses & Not yet covered. \\
B1.5.7.3.10 & Viral diseases not classified in the system: dengue hemorrhagic fever and chikungunya fever; varicella; zoster; herpes simplex infection; viral exanthems (measles, rubella, erythema infectiosum, exanthema subitum, and hand, foot, and mouth disease); mumps; infectious mononucleosis; and cytomegaloviral infection & Not yet covered: these diseases are not discussed clinically. \\
\bottomrule
\end{tabular}
\end{center}


\section{Introduction to Medical Virology}

\marginnote{\begin{kaobox}[frametitle={Origins of selected virus names},backgroundcolor=RoyalBlue!12!White,frametitlebackgroundcolor=RoyalBlue!25!White]
\label{box:virus-naming}
\begin{itemize}
  \item \textit{Picornaviridae}\index{Picornaviridae} are small RNA viruses; the name combines ``pico'' (small) and RNA.
  \item \textit{Togaviridae}\index{Togaviridae} have a membrane envelope; the name derives from Latin \textit{toga}, meaning ``cloak.''
  \item Reoviruses\index{reovirus} are named from ``respiratory, enteric, orphan''; they have segmented double-stranded RNA genomes and do not use reverse transcription.
  \item Poxviruses\index{poxvirus} are named for pock-like lesions; the family includes variola virus\index{variola virus}, which causes smallpox.
  \item Adenoviruses\index{adenovirus} were first isolated from adenoid tissue; the family is named \textit{Adenoviridae}\index{Adenoviridae}.
\end{itemize}
\end{kaobox}}

\textbf{Viruses}\index{virus} may be defined as microscopic, acellular, \textbf{obligate intracellular infectious agents}\index{obligate intracellular agents} that replicate only within host cells. \textbf{Viral replication}\index{viral replication} depends on host-cell resources and viral components to produce progeny\index{progeny}, which are referred to as virions\index{virion} (they are the extracellular form of a virus).\footnote{They do not reproduce independently or by binary fission, and hence are not typically considered living organisms.} The majority of viruses cause disease, but viruses vary in their interactions with and evasion of host defenses. Medically important viruses often share characteristics, including (1) the ability to infect susceptible cells\index{succeptible cells} and produce viral particles; (2) interactions with host immune defenses, including mechanisms of \textbf{immune evasion}\index{immune evasion}; and (3) \textbf{transmission between hosts}\index{viral transmission} by defined means, often through protective barriers (e.g., contact, respiratory droplets or aerosols, fecal--oral transmission, vectors, or other routes).

Viruses consist of genetic material (either \textbf{DNA}\index{viral DNA}, deoxyribonucleic acid, or \textbf{RNA}\index{viral RNA}, ribonucleic acid) enclosed in a \textbf{protein capsid}\index{capsid}.\footnote{One subunit of a protein capsid is referred to as the \textbf{capsomere}\index{capsomere}.} Some viruses possess a lipid envelope\index{viral envelope} surrounding their capsid. These viruses are referred to as \textbf{enveloped viruses}. Alternately, non-enveloped viruses are \textbf{naked viruses}\index{naked virus}. Both enveloped and naked viruses may present spikes on the surface of their virion: for naked viruses, these spikes present on the protein capsid, while for enveloped viruses, they present on the outermost envelope.

Classification of viruses utilizes factors such as genome type and replication strategy (following the \textbf{Baltimore classification}), capsid symmetry, and the presence or absence of an envelope. Notably, formal classification systems are also used in practice. Virus names may reflect characteristics such as size, morphology, target infection site, patient symptoms. Refer to Box~\ref{box:virus-naming} for examples of nomenclature for selected common viruses.

\section{The Classification of Viruses}

Four commonly used descriptors for virus classification are (1) genetic composition; (2) capsid symmetry\index{capsid symmetry}\footnote{Surprisingly not much different from a `Dynamite' candy with its soft chocolate core; you can tell I am hungry.}; (3) the presence or absence of a lipid envelope; and (4) the dimensions of the virion\index{virion} or capsid. These features help describe and compare viruses, but formal taxonomy also considers genetic and evolutionary relationships.

Capsid morphology provides a useful structural description. Common forms include helical and icosahedral capsids, while some virions have more complex structures. Envelope status is a separate feature: enveloped virions also have a capsid beneath their lipid membrane.


\begin{center}
  \includegraphics[width=\textwidth]{\booknlroot appendices/appendix 3: Medical Microbiology and Immunology/images/c2-capsid-shapes.png}
  \captionof{figure}{Three common viral capsid morphologies: helical, icosahedral, and complex. The bacteriophage shown has an icosahedral head and a protein tail with tail fibers. Reproduced from CK-12 Foundation, \textit{The Living Environment}, section 14.2~\cite{CK12CapsidShapes}.}
  \labfig{c2-capsid-shapes}
\end{center}

We first focus on classification based on the outer shell of a virus. The outer layer protects the viral genome and assists in delivery to target cells through \textbf{viral attachment proteins (VAPs)}\index{viral attachment protein}.\footnote{Antibodies that block VAPs can prevent viral entry and neutralize some viruses.} The viral genome is enclosed in a protein coat (capsid); enveloped viruses also have a lipid envelope surrounding the capsid.

\section{Non-enveloped Viruses}

Non-enveloped viruses are generally more resistant than enveloped viruses to detergents, drying, acids, and bile due to their rigid protein capsids. This allows many of them to be transmitted by the fecal--oral route\index{fecal-oral transmission}.

\section{Enveloped Viruses}


On the other hand, \textbf{enveloped viruses}\index{enveloped virus}, because of their lipid envelopes, are more susceptible to drying, detergents, and solvents. Therefore, they are typically transmitted through fluids (including bodily secretions, blood, and tissue), and many do not survive in the gastrointestinal tract, although this varies among viruses.

\section{The Baltimore Classification}

The Baltimore classification for viruses classifies viruses based on their genetic material.

\begin{center}
  \includegraphics[width=\textwidth]{\booknlroot appendices/appendix 3: Medical Microbiology and Immunology/images/c2-baltimore-classification.png}
  \captionof{figure}{The seven groups of the Baltimore classification, organized by viral genome type and the pathway used to produce mRNA and viral proteins. Reproduced from ViralZone~\cite{ViralZoneBaltimore}.}
  \labfig{c2-baltimore-classification}
\end{center}

\section{Viral Infection}

\subsection{Staging of Viral Infections}
\label{stagingvirus}
Viral infections are either productive\index{productive infection} or abortive\index{abortive infection} (failed infections), depending on both factors determined by the host cell and virus. For an infection to succeed, the cell must be susceptible and permissive\footnote{After all, it should come as no surprise that a virus must first be able to infect the cell, as well as replicate within it to infect the host. Otherwise, an infection will become abortive, which means that new viral progeny is produced.}. A susceptible cell\index{succeptible cell (to virus)} is a cell which expresses receptors that may bind viral proteins and allow entry, while a permissive cell\index{permissive cell (to virus)} contains the intracellular substances required for a virus to proliferate.

A successful infection may be classified further based on the duration. Infections may either be acute or chronic, and chronic infections may be further classified based on the presence or absence of a detectable viral load. A latent infection\index{latent infection} presents without active viral load, but with the ability to re-infect if the host immunity drops. On the other hand, chronic infections\index{chronic infection} present with a persistent (and typically reduced viral load compared to the acute phase of infection). Therefore, chronic infections often present as asymptomatic. Refer to \ref{c2-viral_stage} for more details.

\begin{center}
  \includegraphics[width=\textwidth]{\booknlroot appendices/appendix 3: Medical Microbiology and Immunology/images/c3-infect-stages}
  \captionof{figure}{General patterns of infection. Relative virus particle production is plotted as a function of time after infection
(blue line). The time during which symptoms are present is indicated by the orange shaded area. Reproduced from~\cite{flint2015principles}.\textbf{ find a higher res version.}}
  \labfig{c2-viral_stage}
\end{center}

\subsection{Pathogenesis of Viral Infection}

The human body typically has a number of innate and adaptive defenses against infections (covered in:~\ref{immune_respond_chapter}). These defenses serve to disrupt the infective process of pathogens. Infections occur when these mechanisms are weakened or bypassed (such as immunocompromised or breaches of the cutaneous barrier). The distribution of a viral infection can best be described as either localized or systemic. Localized infections\index{localized infection} refer to colonization at and around the site of infection, while
disseminated infections\index{disseminated infection} refer to those which have spread to distant sites.\footnote{This typically occurs through viremia\index{viremia} which refers to the presence of a virus within the bloodstream, which may occur transiently.} The route of entry for viruses governs (to some extent) the patterns of infection and subsequent pathology that are expected. Examples of disseminated spread include: Invasion of sub-epithelial and lymphatic tissue, vascular endothelial invasion, persistent viremia and macrophage or nerve infection/spread.

The major barrier mechanisms of the immune system against viruses are:\footnote{refer to~\ref{immune_respond_chapter} for details}
\begin{itemize}
  \item Anatomical and chemical barriers -- barriers that inactivate viruses or prevent penetration of viruses into susceptible cells. This includes physical barriers and secreted substances such as mucus, saliva, and defensins, as well as physical barriers such as skin (the majority of infections stop at this step).
  \item Intrinsic barriers -- barriers which prevent viruses from taking root in host cells. This includes cellular mechanisms such as apoptosis, autophagy and interferons.
  \item Innate immunity -- the immune system's immediate and nonspecific response to pathogens through recognition of pathogenic factors (often pathogen associated molecular patterns). This involves phagocytes, NK cells, antigen presenting cells (APCs) and neutrophils.
  \item Acquired immunity -- the immune system's customized response for each pathogen presented by APCs. This involves the production of antibodies by immune cells and targeted attack of the pathogen.
\end{itemize}

For a viral infection to take root, the virus must: (1) have a mechanism of bypassing host immunity; (2) be present in an amount exceeding the innoculum dose; (3) cells must both be susceptible and permissive (see figure~\ref{stagingvirus}), and (4) host immunity must be insufficient to defend against the infection. If one or more of these factors are not present, viral infection will not occur.

The viral (infective) titer of a virus\index{viral titer} is a measure used to establish infectivity and it is defined as the minimum inoculum of a virus required to establish an infection. Various experimental assays may be used to make this determination.\footnote{Hence, the specific value varies by assay platform.} Example platforms include: tissue culture (TCID\textsubscript{50} and PFU) and animal culture (ID\textsubscript{50} and LD\textsubscript{50}).

% types of viral infections

\subsection{Methods of Viral Entry}

\subsection{Pathogenesis of Viral Infections}

\subsection{Epidemiology \& Mechanisms of Viral Transmission}


\end{document}